\honest{Sneaking in Connotations}{Eliezer Yudkowsky}{2008}

Yesterday we saw that in Japan, blood types have taken the place of astrology. So suppose we invent a word, ``wiggin'', and define it to mean people with green eyes and black hair.

Danny sees a man with green eyes and black hair walk into a restaurant. ``Bloody wiggins,'' he says. ``Commit all sorts of crimes, they do.'' His sister Erda asks whether he has seen the man commit any. Danny has no need to: the dictionary says a wiggin is a person with green eyes and black hair, and the man has both. Just watch, he says; the man will reach for the ketchup.

Words carry the mind from what it sees to what it infers. A dictionary lists the features by which you recognize a thing, not ``the ten thousand connotations'' you infer from it. I offer two guesses why, each with ``perhaps''. So Danny's dictionary lists green eyes and black hair, and says nothing of crime or ketchup. How did those get attached to the word? Maybe someone made things up and wrote bestselling books about it. \nb{That fits the real case from the previous post. The Japanese belief in blood-type personalities is traced to books published in the 1970s by Masahiko Nomi, a journalist with no medical training.}

A definition that included the inferred property could not be applied until you had seen the property. So ``Whenever someone pulls a dictionary, they're generally trying to sneak in a connotation.'' If ``wiggin'' only means green-eyed and black-haired, why not call those people ``green-eyed black-haired people''? \nb{One short word is easier to use than four. Yudkowsky had written six days earlier that ``words we use frequently should be short''.} And if you want to know whether he reaches for the ketchup, why not ask that?

Arguing the real question takes work, ``And people are lazy.'' I give no evidence for this. The real reason people care whether someone is a wiggin is the feeling that comes with the word.

My test: imagine Danny arguing from the dictionary that the man is a wiggin and ``therefore, he's got black hair.'' It has no triumphant ring. An argument that stayed inside the definition would ``either say nothing new, or beg the question,'' so only a smuggled connotation makes anyone feel they can score a point that way. \nb{That holds for one-step arguments like Danny's. A long deduction from definitions, as in mathematics, can tell a hearer something new, and the next post exempts mathematics.}
